% Edited conversion of source pp. 17--18.
\ReportHeading
  {Моделювання структурних залишкових напружень у порожнистому сталевому циліндрі, зумовлених технологічним нагріванням}
  {О.~Р. Гачкевич, Т.~В. Козакевич, Т.~Волчаньскі}
  {Інститут прикладних проблем механіки і математики ім. Я.~С. Підстригача НАН України, Львів, Україна; Опольська політехніка та Освітньо-навчальне об’єднання у Глушині, Польща}
  {}

Для прогнозування механічних властивостей конструктивних елементів і рівня залишкових напружень за різних способів нагрівання та охолодження важливо розробити методику дослідження й оптимізації фазового та напруженого станів сталевих циліндричних елементів як під час термообробки, так і після її завершення.

До розглядуваного класу задач адаптовано відому математичну модель фазового складу та зумовлених ним залишкових напружень у маловуглецевих низьколегованих сталевих тілах за монотонного охолодження з аустенізованого стану. На цій основі сформульовано розрахункову схему задачі механіки з урахуванням структурних перетворень. Її реалізують послідовно: розв’язують задачу теплопровідності, визначають метастабільний фазовий склад і розраховують залишкові напруження.

Задачу про напружений стан порожнистого циліндра розв’язано методом скінченних елементів із використанням експериментально встановленої залежності коефіцієнта теплообміну для конкретного охолоджувального пристрою.

На основі запропонованої методики досліджено вплив умов охолодження зовнішньої та внутрішньої поверхонь на фазовий склад і залишковий напружений стан вільного від зовнішнього силового навантаження порожнистого циліндра зі сталі 17Г1С.

Результати показують, що ділянки, які охолоджуються інтенсивніше, містять більшу частку гартувальних складових --- мартенситу та бейніту. У цих ділянках виникають локалізовані стискальні напруження, що добре узгоджується з експериментальними даними.

Розроблені модель і методика розрахунку фазового складу можуть бути використані для аналізу й оптимізації ширшого класу режимів охолодження сталевих тіл з метою цілеспрямованого керування їх структурою, залишковим напруженим станом і механічними властивостями.

Отримані числові оцінки фазового складу та залишкових напружень корелюють з емпіричними даними: частка гартувальних складових, зокрема мартенситу, у приповерхневих ділянках загартованих тіл більша, ніж у центральних. Оскільки мартенсит має меншу густину, ніж рівноважна феритно-перлітна структура, у збагачених ним ділянках виникають стискальні напруження. Наведено числові оцінки цих напружень для різних режимів охолодження зовнішньої та внутрішньої поверхонь порожнистого сталевого циліндра.

\renewcommand{\refname}{Список використаних джерел}
\begin{thebibliography}{9}
\bibitem{residual:gachkevych2023}
О.~Гачкевич, Т.~Козакевич, Р.~Кушнір,
``Вибрані математичні проблеми термомеханіки маловуглецевих низьколегованих сталевих пластин при нагріванні рухомими джерелами тепла з урахуванням структурних залишкових деформацій,''
\emph{Праці Наукового товариства ім. Шевченка}, т.~LXXIV, с.~81--98, 2023.
\end{thebibliography}
